% De Lege Agraria II (de-lege-agraria-2) --- English translation
% Translated by Claude (Anthropic) from the Latin (Perseus).
% Style guide: STYLE.md.

\ciceroSection{1} It is established by custom, citizens, and by the institution of our ancestors, that those who through your kindness have attained the right of family portrait-busts should hold their first public address to link the gratitude they owe your kindness with the praise of their own forebears. In which speech a few are sometimes found worthy of their ancestors' place; but most achieve only this --- to make it look as though so much was owed to their ancestors that even something was left over to be paid out to their descendants. To me, citizens, it is not granted to speak before you about my ancestors --- not that they were not such men as you see those of us were who were begotten of their blood and trained by their teachings, but because they had no part in popular renown or in the light of public office.

\ciceroSection{2} About myself, however, I fear that to speak before you would be arrogant, to keep silent ungrateful. For both to recall by what efforts I have attained this dignity is exceedingly hard, and to be silent about your great kindnesses I cannot in any way. Therefore I shall observe a fixed measure and moderation in speaking: I shall recall what I have received from you; why I am worthy of your highest honour and your singular judgment, I shall myself say with restraint, if it must be said, supposing that you yourselves who passed the verdict will think the same.

\ciceroSection{3} You made me, after a very long interval reaching nearly to the limit of our memory, the first ``new man'' to be made consul; and that place which the nobility held fortified by garrisons and walled around by every device, you tore open under my leadership and willed it to lie open to virtue for the future. And you made me consul not only --- which is in itself most splendid --- but you made me so in a way in which few nobles in this city have been made consuls, and no new man before me. For if you will recall the new men, you will find that those who were made consul without a defeat were made so by long labour and some opportunity, when they had stood many years after their praetorship and rather later than their age and the laws would have allowed; while those who stood in their year did not become consul without a defeat. I am the only one of all the new men we can remember who stood for the consulship as soon as it was permitted, was made consul as soon as I stood --- so that your honour was sought to the day of my own time, not snatched up at the chance of another's candidacy, not extorted by long entreaties, but won by dignity.

\ciceroSection{4} That is the most splendid thing I mentioned just now, citizens --- that with this honour you adorned me, the first new man for many years, at my first standing, in my year. But still, more magnificent and more ornate than that, nothing can be: that at my election you cast not the silent ballot, that vindicator of secret liberty, but you carried before you a living voice as the index of your goodwill and zeal toward me. And so it was not the final counting of votes, but those first thronging assemblies of yours, not the single voices of the heralds, but the one voice of the entire Roman people, that proclaimed me consul.

\ciceroSection{5} This kindness of yours, citizens, so distinguished and so singular, I count great indeed for the joy and gladness of my mind, but greater still for my care and anxiety. For there are turning in my mind many heavy thoughts that grant me no part of either day's or night's rest: first, the keeping of my consulship --- a charge that is hard and great for everyone, but for me beyond the rest, since for me an error finds no pardon, a thing rightly done finds slim praise wrung out of unwilling givers; faithful counsel is not shown to me in doubt, certain support is not offered me in toil from the nobility.

\ciceroSection{6} If I alone were brought into some peril, I should bear it, citizens, with a more even mind. But it seems to me certain men, if in any matter they shall judge me to have slipped --- not in counsel only, but even by accident --- will censure all of you who set me before the nobility. For my part, citizens, I judge that I must endure all things rather than fail to conduct my consulship in such a way that in all my actions and counsels your action and counsel about me may be praised. There is added to me also that highest labour and most difficult method of holding the consulship, that I have decreed I must not use the same law and condition as previous consuls, who shunned approach to this place and the sight of you in part with great care, in part did not pursue it with vigour. But I, not only in this place, where it is most easily said, but in the Senate itself, where there seemed no place for such a voice, declared in my very first speech on the Kalends of January that I would be a consul friendly to the people.

\ciceroSection{7} For I cannot in any way bring it about that, when I see myself made consul not by the zeal of powerful men nor by the surpassing favour of a few, but by the judgment of the entire Roman people, made consul in a way that I should be set far before the most noble men, I should not seem to be a friend of the people both in this magistracy and in the whole of my life. But for the force and interpretation of this word I have great need of your wisdom. For a great error is in circulation because of the treacherous pretences of certain men who, when they fight against and obstruct the people's not only advantages but also safety, wish by their oratory to seem friends of the people.

\ciceroSection{8} What kind of commonwealth I received on the Kalends of January, citizens, I see well: full of anxiety, full of fear, in which there was no evil, no calamity, that the loyal did not dread and the wicked await; all turbulent designs against the present condition of the commonwealth and against your peace were said to be in part being undertaken, in part already undertaken when we were consuls-elect. Faith had been driven from the Forum, not by any single blow of fresh disaster, but by suspicion and the disturbance of the courts, the weakening of cases already judged. New dominations, extraordinary not commands but kingships, were thought to be the object of pursuit.

\ciceroSection{9} When I not only suspected this but plainly saw it --- for it was not being done in secret --- I said in the Senate that in this magistracy I would be a consul friendly to the people. For what is so dear to the people as peace? in which not only those whom nature has endowed with sense, but even roofs and fields, seem to me to rejoice. What so dear to the people as liberty? which you see sought not only by men but even by beasts, and set before all things. What so dear to the people as quiet? which is so pleasant a thing that you and your ancestors and every bravest man have thought the greatest labours should be undertaken so that one might at last be in quiet, especially in command and dignity. Indeed, on this account also we owe to our ancestors a particular praise and gratitude --- that by their labour it has been brought about that we may be in quiet without fear. How then can I not be a friend of the people, when I see all these things, citizens --- foreign peace, the liberty proper to your race and name, domestic quiet, in short all the things dear and ample to you --- entrusted to the keeping and as it were the patronage of my consulship?

\ciceroSection{10} For neither, citizens, ought it to seem to you pleasant or popular if some bounty has been promulgated, the kind of thing that can be paraded in words but in truth cannot be done unless the treasury is drained dry; nor are these to be reckoned popular: disturbances of the courts, the undoing of cases already judged, the restoration of the condemned --- which in stricken states, when affairs are already beyond hope, are wont to be the last terminations of disasters. Nor, if any men promise lands to the Roman people, if they are pushing one thing in secret while they parade another in hope and the colour of pretence, are they to be thought friends of the people. For I will speak truly, citizens: the very kind of agrarian law I cannot blame. For there comes to my mind that two most distinguished men, of the highest talent, most loving toward the Roman commons, Tiberius and Gaius Gracchus, settled the commons on public lands which had been previously held by private men. I am not, however, the kind of consul who, like very many, judges it impious to praise the Gracchi, by whose counsels, wisdom, and laws I see many parts of the commonwealth set in order.

\ciceroSection{11} Therefore, when at the start, just consul-designate, it was being announced to me that the tribunes-designate were drafting an agrarian bill, I was eager to learn what they had in mind. For I judged that, since we had to hold our offices in the same year, there ought to be some partnership between us in the well-administering of the commonwealth.

\ciceroSection{12} When I tried to insinuate myself in friendly fashion into their conversations and to lend my services, I was kept in the dark, shut out; and when I made it clear that, if a law were to seem to me useful to the Roman commons, I would be its sponsor and helper, they nevertheless rejected this generosity of mine and denied that I could be brought to approve any bounty. I made an end of offering myself, lest perhaps my eagerness should seem either treacherous or shameless. Meanwhile they did not cease to meet privately among themselves, to bring in certain private persons, to add the night and solitude to their secret gatherings. By which things how great a fear we were in, you can easily make out by conjecture from the anxiety in which you were yourselves at that time.

\ciceroSection{13} The tribunes of the plebs at last enter on their office. The public address of Publius Rullus is awaited, since he was both the chief of the agrarian bill and was bearing himself more truculently than the rest. When designate already he was practising another expression, another tone of voice, another walk, with shabbier clothing, his body unkempt and shaggy, his hair longer than before and his beard greater, so that by his eyes and looks he seemed to be announcing to all his tribunician violence and threatening the commonwealth. I was waiting for the man's law and his speech: at first no law is published, he orders an assembly to be called for the day before the Ides. There is a great rush of expectation. He unfolds quite a long speech, in very fine words. There was one thing that seemed to me defective: that out of so vast a throng no one could be found who was able to make out what he was saying. Whether he did this for the sake of stratagem, or whether he delights in this kind of eloquence, I do not know. Yet those who had stood somewhat sharper-witted in the assembly conjectured that he had wished to say I-know-not-what about an agrarian law. At length, when I was already designate, the law is published. At my order, several scribes hurry up at the same moment and bring me the law transcribed.

\ciceroSection{14} On every count I can assure you, citizens, that I came to read and study the law in this spirit: that, if I should perceive it to be fitted to your interests and useful, I should be its sponsor and helper. For neither does the consulship by nature, by separation, or by any inherent hatred have undertaken some kind of war with the tribuneship --- because very often loyal and brave consuls have stood up to seditious and wicked tribunes of the plebs, and because tribunician violence has sometimes resisted consular caprice. It is not the unlikeness of the powers that causes the dispute, but the disjunction of minds.

\ciceroSection{15} And so I took the law into my hands in this spirit: that I wished it to be suited to your interests, and of such a kind that a consul popular in deed, not in oratory, could honourably and gladly defend it. And from the first clause of the law to the last I find, citizens, that nothing else has been thought, nothing else undertaken, nothing else done, except this: that ten kings should be set up over the treasury, the revenues, all the provinces, the entire commonwealth, the kingdoms, the free peoples, in short the lords of the world --- under the pretence and the name of an agrarian law. So I assert, citizens, that by this fine and popular agrarian law nothing is given to you, all things are bestowed on certain men, lands are paraded before the Roman people, even your liberty is torn away, the fortunes of private men are increased, those of the state are drained, and finally --- which is most outrageous --- through a tribune of the plebs, whom our ancestors wished to be the protector of liberty and its guardian, kings are set up in the state.

\ciceroSection{16} When I shall have set these things out, citizens, if they shall seem to you false, I shall follow your authority, I shall change my own opinion. But if you perceive that snares are being set against your liberty under the pretence of a bounty, do not hesitate to defend, with no labour of your own, with the consul as your helper, the liberty won by your ancestors at vast cost of sweat and blood and handed down to you. The first clause of the agrarian law is one by which, as they think, you are tested gently in the matter of your liberty's diminution, to see how you can bear it. For he orders that the tribune of the plebs who has carried the law create the ten commissioners through seventeen tribes, so that whomever nine tribes have made, that man shall be a decemvir.

\ciceroSection{17} Here I ask why he took the beginning of his measures and his laws from this point --- that the Roman people be deprived of its vote. So often by agrarian laws have commissioners been set up: triumvirs, quinquevirs, decemvirs; I ask of this popular tribune of the plebs whether they were ever created except through the thirty-five tribes. For when it is fitting that all powers, commands, and commissions should proceed from the entire Roman people, certainly those most of all that are set up for some advantage and benefit of the people, in which both all together pick out the man whom they think will most consult the Roman people, and each one for himself, by his zeal and his vote, can pave a way to obtain the favour. To this tribune of the plebs above all it has come to deprive the entire Roman people of its votes, and to call a few tribes, not on a fixed condition of right but by the favour of a chance lot, to the exercise of liberty.

\ciceroSection{18} ``Likewise,'' he says, ``and in the same way,'' in another clause, ``as in the elections of the Pontifex Maximus.'' He did not even see this, that our ancestors were so popular that, where it was not religiously lawful for a man to be created by the people on account of the religious nature of the rites, in this case nevertheless on account of the magnitude of the priesthood they wished the people to be petitioned. And the same Gnaeus Domitius, tribune of the plebs, a most distinguished man, carried the same measure as to the other priesthoods, that, since the people could not from religious scruple appoint priests, a smaller part of the people should be summoned, and whoever was made by that part should be co-opted by the college.

\ciceroSection{19} See what the difference is between Gnaeus Domitius, tribune of the plebs, a most noble man, and Publius Rullus, who has tested, I suppose, your patience by saying that he is a noble. Domitius, in a thing that could not be done by the people on account of the religious rites, contrived by reason that the part of the people, as far as he could, as far as it was lawful, as far as it was permitted, should have its share. This man, in a thing that always was the people's own, that no one ever lessened, no one ever altered --- that those who were going to assign lands to the people should first receive the favour from the people before they gave it --- has tried to wrench from your hands and snatch away wholly. The other gave to the people in some manner what could in no way be given; this one tries by some method to wrench away what can in no way be taken.

\ciceroSection{20} One will ask, in such injury and such impudence, what end he had in view. He had no lack of design: faith toward the Roman commons, citizens, and fairness toward you and toward your liberty, he sorely lacked. For he orders him who carried the law to hold the elections for the creation of the ten commissioners. I will say it more plainly: Rullus, a man neither greedy nor grasping, orders Rullus to hold the elections. I do not yet object; I see others have done so. The thing no one ever did --- to deal with a smaller part of the people --- see whither it tends. He will hold the elections, he will choose to declare elected those for whom royal power is sought by this law. The whole people he himself does not trust, nor do those authors of these designs think they can rightly be entrusted with it.

\ciceroSection{21} The same Rullus will draw lots for the tribes. The lucky man will draw out the tribes he wants. And those whom nine tribes drawn out by the same Rullus shall have made decemvirs, these we shall have, as I shall presently show, as masters in all things. And these, that they may seem grateful and mindful of the favour, will confess that they owe something to certain men of those nine tribes; but for the remaining twenty-six tribes there will be nothing they will not think they can rightly refuse. What men, then, does he wish made decemvirs? Himself first. By what right? For there are old laws --- and not consular laws, if you think this counts for anything, but tribunician laws, very pleasing and dear to you and your ancestors: there is the Lex Licinia and another, the Lex Aebutia, which exempt not only the man who has carried a measure for some commission and power, but also his colleagues, kinsmen and connections, from having that power or commission entrusted to them.

\ciceroSection{22} For if you are looking out for the people, remove yourself from the suspicion of any advantage to yourself, give proof that you are looking only for the people's profit and gain, allow the power to come to others, the gratitude for the favour to yourself. For this thing is scarcely fitting for a free people, scarcely fitting for your spirit and magnificence. Who carried the law? Rullus. Who barred the larger part of the people from the votes? Rullus. Who presided at the elections, who summoned the tribes he wished, with no overseer of the lots, who created the decemvirs he wished? The same Rullus. Whom did he proclaim chief? Rullus. By Hercules, I scarcely think he could prove this acceptable to his slaves --- not to mention to you, masters of all nations. The best laws, then, will be done away by this law without exception; the same man will seek for himself by his own law a commission; the same man with the larger part of the people robbed of votes will hold elections; whomever he wishes, including himself, he will declare elected; and, of course, he will not refuse his colleagues, the co-signers of the agrarian law, to whom the first place in the index and in the heading of the law has been granted to him; the rest of the fruits of all things which lie set in the hope of this law are kept by joint guarantee and on equal share.

\ciceroSection{23} But see the man's diligence --- if you suppose either that Rullus thought of this, or that it could come into Rullus's mind. Those who contrived these things saw that, if you had been given the power of choosing from the entire people, in any matter where faith, integrity, virtue, authority were sought, you would without hesitation entrust it to Gnaeus Pompey as your chief. For the one man whom out of all you would pick to set in command of all the world's wars on land and sea, certainly in the appointing of decemvirs, whether faith were considered or honour, both could most safely be entrusted to him and most justly adorn him.

\ciceroSection{24} And so this law makes no exception of youth, of any legal impediment, of office, of any magistrate hindered by other duties or laws; even an accused man is not excepted from being eligible to be made decemvir; --- Gnaeus Pompey alone is excepted, lest with Publius Rullus --- I say nothing of the others --- he be made decemvir. For he orders him to declare himself in person, which has never been required in any other law, not even in those magistracies which have a fixed order of standing, lest, if the law were accepted, you should add him as a colleague to Rullus, a guardian and avenger of his greed. Here, since I see that you are stirred by the man's worth and by the law's insult, I shall renew what I said at the start --- a kingship is being prepared, your liberty by this law is being utterly taken away.

\ciceroSection{25} Or did you think otherwise? when a few men cast eyes of greed upon all that is yours, did you not think they would in particular act so as to drive Gnaeus Pompey from every guardianship of your liberty, from every power, commission, and patronage of your interests? They saw and they see that, if through your imprudence and my negligence you accept a law you do not know, it will come about that, when later you have learned the snares set for you, when you make the decemvirs, then you must think Pompey's protection should be opposed to all the law's flaws and crimes. And will it not be a small proof to you that domination and power over all things is being sought by certain men, when you see that the man whom they perceive to be the guardian of your liberty is being made to have no part in the dignity?

\ciceroSection{26} Now learn what power is given to the decemvirs, and how great. First, by a curiate law he equips them. This is unheard of and quite new --- that magistracy be given by a curiate law to those who have not first been given it by any election. He orders that law to be carried by that praetor of the Roman people who has been made first. But how? That those should hold the decemvirate whom the plebs has appointed. He has forgotten that none are appointed by the plebs. And does this man bind the world with new laws who does not remember in his third clause what he has written in his second? Here it is plain what right you have received from your ancestors, what is left to you by this tribune of the plebs. Our ancestors wished you to give your judgment twice on each magistracy. For when the centuriate law was carried for censors, when the curiate for the other patrician magistrates, then a second judgment was given on the same men, so that there might be the power of recalling the favour, if the people came to repent.

\ciceroSection{27} Now, citizens, you keep those first elections, the centuriate and the tribute; the curiate are kept only for the sake of the auspices. But this tribune of the plebs, because he saw that no one could hold a power without the order of the people or the plebs, has confirmed it by curiate elections which you do not enter, has done away with the tribute elections which were yours. So when our ancestors wished you to judge concerning each magistracy by two elections, this popular man has not left even one election in the people's power.

\ciceroSection{28} But see the man's scrupulousness and diligence. He saw and clearly perceived that without a curiate law the decemvirs could not have power, since they were appointed by nine tribes; he orders a curiate law to be carried about them; he commands the praetor. How absurdly he does it does not concern me. He orders the praetor who has been made first to carry the curiate law; if he cannot carry it, the one who has been made last --- so that he seems either to have toyed in such great matters or assuredly to have had I-know-not-what in view. But this thing, which is either so perverse as to be ridiculous, or so malicious as to be obscure, let us pass over; let us return to the man's scruples. He sees that without a curiate law nothing can be done by the decemvirs.

\ciceroSection{29} What then, if it shall not be carried? Mark the cunning. ``Then,'' he says, ``those decemvirs shall be of the same right as those who have been created under the best law.'' If this can be done in this state which far surpasses other states in the right of liberty --- that anyone obtain command or power without any election --- what is the use of ordering, in the third clause, a curiate law to be carried, when in the fourth you allow that without a curiate law they shall have the same right as if they had been created by the people under the best law? Kings are being set up, citizens, not decemvirs. And so they take their birth from these beginnings and foundations such that, not only when they shall have begun to discharge the magistracy, but even when they are being set up, every right, power, and liberty of yours is taken away.

\ciceroSection{30} See now how diligently he keeps the right of the tribunician power. When consuls have carried curiate laws, the tribunes of the plebs have often vetoed them --- nor do we complain that this is the tribunes' power; only, if anyone has used that power rashly, we count him a madman --- but here this tribune of the plebs takes away from the curiate law that the praetor will carry the power of veto. And while this is to be condemned in him because by means of a tribune of the plebs the tribunician power is diminished, it is also to be laughed at, because while a consul, if he has not got a curiate law, may not touch military business, yet this man, in the very case where he forbids veto, sets up a power that, even if a veto is interposed, will still be the same as if the law had been passed --- so that I cannot see why he either forbids veto or thinks anyone will veto, since the veto will only signify the vetoer's folly and not impede the matter.

\ciceroSection{31} Let the decemvirs, then, be set up neither by true elections --- that is, by the votes of the people --- nor by those shadowy elections, kept up for show and for the maintenance of antiquity, with the thirty lictors for the auspices' sake. See now what greater ornaments he confers upon those who have received from you no power, than upon all of us upon whom you have bestowed the most ample powers. He orders the decemvirs to have the auspices for planting colonies, and the keepers of the sacred chickens, ``with the same right,'' he says, ``as the triumvirs had under the Lex Sempronia.'' Do you dare even, Rullus, to make mention of the Lex Sempronia? does that very law not remind you that those triumvirs were created by the votes of the thirty-five tribes? And, since you are very far removed from Tiberius Gracchus's fairness and modesty, do you suppose that what was done on a most unlike footing ought to be of the same right?

\ciceroSection{32} Besides, he gives a power praetorian in name, in fact regal; he limits it to five years; he makes it perpetual, for he secures it with such resources and forces that against the unwilling it cannot in any way be taken away. Then he equips them with attendants, scribes, copyists, heralds, architects, and besides with mules, tents, saddle-blankets, baggage. He drains the expense from the treasury, replenishes it from the allies. He sets up two hundred surveyors of equestrian rank, twenty bodyguards apiece, the same as ministers and satellites of the power. So far you have only the form, citizens, only the appearance of tyrants; you see the trappings of power, not yet the power itself. For someone may perhaps say: ``what hurt do these things do me, the scribe, the lictor, the herald, the keeper of the sacred chickens?'' All these things are of such a kind, citizens, that the man who holds them without your votes seems either an unbearable king or a private madman.

\ciceroSection{33} Look at how great a power is granted; you will say it is not the insanity of private men but the unrestraint of kings. First, infinite power is given of raising innumerable money by the alienation, not the enjoyment, of your revenues. Then, over the world and over all nations, judicial cognizance is given without a council, punishment without appeal, censure without recourse to help.

\ciceroSection{34} For five years they will be able to pass judgment on consuls or on the very tribunes of the plebs; meanwhile no one will pass judgment on them. They will be allowed to seek magistracies, not allowed to plead a case in court. They will be able to buy lands from whomever they wish and at whatever price they wish, however dearly. They are allowed to plant new colonies, renew old ones, fill all Italy with their colonies. The supreme power is given of going through every province, of fining free peoples in their lands, of selling kingdoms. When they wish, they may be at Rome; when it suits them, anywhere they please they may roam with the supreme imperium and judgment of all things. Meanwhile let them dissolve public courts, withdraw from councils whomever they wish, decide each one alone of the greatest matters, hand it over to a quaestor, send a surveyor; let what the surveyor reports to the one man by whom he was sent be ratified.

\ciceroSection{35} A word fails me, citizens, when I call this power kingly --- but in fact a certain greater word is needed. For there was never any kingdom that did not contain itself, if not within some right, then at least within fixed regions. But this power is boundless: by the law's permission it embraces all the kingdoms, and your empire which extends most widely, and those things which are partly free from you, partly even unknown to you. They are given, then, first, the power to sell all those things which it was decreed by the Senate should be sold in the consulship of Marcus Tullius and Gnaeus Cornelius and afterwards.

\ciceroSection{36} Why is this so dim and unseeing? Could not all those things which the Senate decreed sold be set down by name in the law? There are two causes for this obscurity, citizens: one of shame --- if there can be any shame in such marked impudence --- the other of crime. For he does not dare to call by name those things the Senate decreed sold; for there are public places of the city, there are little chapels which since the restoration of the tribunician power no one has touched, which our ancestors wished to be partly the city's ornaments, partly refuges in danger. These the decemvirs will sell by tribunician law. To this will be added Mount Gaurus, the willow-groves at Minturnae, even that very saleable Herculanean way of many delights and great cash, and many other things which the Senate decreed sold from the narrowness of the treasury, but which the consuls did not sell because of the odium.

\ciceroSection{37} But these things perhaps from a sense of shame are passed over in silence in the law. The thing more to be believed and more to be feared is this: that to the decemviral audacity a great power is given of corrupting the public records and of forging decrees of the Senate which were never made --- since out of the number of those who were consuls in those years many are dead. Unless perchance it is unfair for us to suspect any boldness on the part of those whose greed seems to find the world too narrow.

\ciceroSection{38} You have one kind of selling, which I see seems great to you. But turn your minds to what follows; you will see this made into a kind of step and approach to the rest. ``Whatever lands, whatever places, buildings.'' What is there besides? Many things in slaves, in cattle, in gold, in silver, in ivory, in clothing, in furniture, in other things. Why should I say more? Did he think it would be invidious if he had named everything? He did not fear odium. What then? He thought it long, and he was afraid he might leave something out; he wrote in ``or anything else'' --- by which brevity you see that nothing is excepted. So whatever there is outside Italy that has been made public property of the Roman people in the consulship of Lucius Sulla and Quintus Pompeius or after, that he orders the decemvirs to sell.

\ciceroSection{39} By this clause, citizens, I say that all races, nations, provinces, and kingdoms have been handed over and bestowed on the decemvirs' authority, judgment, and power. First I ask this: is there any place anywhere that the decemvirs cannot say has been made public property of the Roman people? For when the same man who has declared so can also pass judgment, what is there that he cannot say, when it is allowed to him to pass the judgment as well? It will suit them to say that Pergamum, Smyrna, Tralles, Ephesus, Miletus, Cyzicum, in short the whole of Asia which has been recovered since the consulship of Lucius Sulla and Quintus Pompeius, has been made the Roman people's;

\ciceroSection{40} will speech be wanting for the discussion of that matter, or, when the same man both discusses and judges, can he not be brought to give a false judgment? or, if he refuses to condemn Asia, will he not appraise the terror of condemnation and his threats at as much as he wishes? What of this, that one cannot in any way speak against, what has now been settled and judged by us, since we have already entered upon the inheritance --- the kingdom of Bithynia, which certainly has been made public property of the Roman people: is there any reason why all the lands, cities, lakes, harbours, the whole of Bithynia in short, will not be sold by the decemvirs? What of Mytilene, which certainly has become yours, citizens, by the law of war and the right of victory, a city famous in the first rank for its nature and site and the layout of its buildings and its beauty, with pleasing and fertile lands? They are included, of course, under the same clause.

\ciceroSection{41} What of Alexandria and all Egypt --- how secretly it lies hid, how it is concealed, how stealthily the whole of it is handed over to the decemvirs! For who of you does not know that this kingdom is said to have been made the Roman people's by the will of King Alexa? Here I, the consul of the Roman people, not only pass no judgment but do not even bring out an opinion. For this great matter seems to me of weight, not only for deciding but even for speaking. I see those who affirm that the will was made; I am aware that the authority of the Senate exists for the entry on the inheritance --- when on Alexa's death we sent ambassadors to Tyre to recover the money he had deposited there.

\ciceroSection{42} I remember that Lucius Philippus often confirmed this in the Senate. Almost all agree that the man who at the present time holds that kingdom is neither of royal birth nor of royal mind. It is said on the other side that there is no will, that the Roman people ought not to seem grasping after every kingdom, that our men will migrate to those parts because of the goodness of the land and the abundance of all things.

\ciceroSection{43} On so great a matter Publius Rullus with the rest of his decemviral colleagues will pass judgment --- and which way will he judge? for either way is so great that in no way can it be conceded or borne. He will wish to be popular; he will adjudge it to the Roman people. So the same man will sell, by his own law, Alexandria, will sell Egypt, will be found judge, arbiter, master --- in short, king --- of a most wealthy and beautiful land. He will not, you say, take so much for himself; he will not grasp it; he will judge that Alexandria belongs to the king, and will judge it away from the Roman people.

\ciceroSection{44} First, why should the decemvirs judge concerning the inheritance of the Roman people, when you have wished the centumviri to judge concerning private inheritances? Then, who will plead the cause of the Roman people? Where will that case be tried? Who are these decemvirs whom we are to see adjudge the kingdom of Alexandria gratis to Ptolemy? But if Alexandria was being sought, why did they not run the same races at this time as they did in the consulship of Lucius Cotta and Lucius Torquatus? Why did they not, openly as before, why not in the same fashion as then, head straight and openly for that region? Or did they who could not hold the kingdom by the etesian winds and a straight course now think they would reach Alexandria by blind darkness and obscurity?

\ciceroSection{45} And consider this in your minds together, citizens. Foreign nations can scarcely bear our ambassadors, men of slight authority, who go on free embassies for the sake of private business. For the very name of empire is heavy, and is dreaded even in a slight person, because it is your name, not their own, that they make use of when they have set out from here. What do you suppose --- when these decemvirs with command, with fasces, with that picked youth of surveyors will roam through the whole world --- with what spirit, what fear, what danger the wretched nations will be?

\ciceroSection{46} There is terror in command; they will endure it. There is expense at his coming; they will bear it. Some service will be commanded; they will not refuse. But this --- how great it is, citizens! --- when that decemvir who has come to some city either expected as a guest or suddenly as a master shall declare that very place to which he has come, that very seat of hospitality into which he has been led, public property of the Roman people! What a calamity to the people if he shall say so --- what a profit to himself if he shall deny it! And the same men who grasp at these things sometimes complain that all lands and all seas have been entrusted to Gnaeus Pompey. As though it were the same thing to entrust many things and to bestow all things, to set a man over labour and trouble or over plunder and gain, to send him to liberate allies or to oppress them! Finally, if some honour is singular, does it make no difference whether the Roman people confers it on whom it will, or whether it is impudently filched from the Roman people by deceit of law?

\ciceroSection{47} You have understood how many things and how great the decemvirs will sell by the law's permission. It is not enough. When they have filled themselves with the blood of allies, of foreign nations, of kings, let them cut the sinews of the Roman people, lay hands on your revenues, break in upon the treasury. For there follows the clause in which he does not even allow --- if perchance money should be lacking, though so much can be received from the previous clauses that it ought not to be lacking --- but plainly, as if this thing would be your salvation, so he forces and orders the decemvirs to sell your revenues by name, citizens.

\ciceroSection{48} Read me out in order, from the law's text, the auction of the Roman people; which by Hercules I think will be a sorrowful and bitter announcement for this very herald.--- ``The auction.'' --- As in his own affairs, so in the commonwealth, he is a luxurious wastrel, who would sell the woods before the vineyards! You have surveyed Italy; pass on to Sicily. --- There is nothing in this province either in the towns or in the lands which our ancestors left to us as our own that he does not order to be sold.

\ciceroSection{49} What our ancestors left you, won by recent victory in the cities and territories of allies as both a bond of peace and a monument of war --- shall you, having received it from them, sell it on this man's motion? Here for a moment I seem to myself to stir your minds, citizens, while I lay open to you the snares which they have, in their own opinion, set deep against Gnaeus Pompey's dignity. And forgive me, I beg, if I appeal to such a man more than once. You, citizens, two years ago, when I was praetor, in this very place imposed on me this role: to defend, with you and so far as I could, the dignity of him in his absence. So far I have done what I could, neither led on by his friendship nor by hope of office and the most ample dignity, which I, though with his goodwill yet in his absence, attained through you.

\ciceroSection{50} For which reason, since I see that this whole law is fitted as a sort of engine for the overthrow of his power, both I shall resist the men's designs and I shall surely accomplish this --- that what I see, you all may not only see but hold fast. He orders to be sold what belonged to the Attalenses, the Phaselites, the Olympeni, and the Aperan, Oroandic, and Gedusan lands. These were made yours by the command and victory of Publius Servilius, that most illustrious man. He adds royal lands of Bithynia which the publicans now enjoy; then the Attalic lands in the Chersonese, in Macedonia, which were of King Philip or Perseus, which were also let out by the censors and yield a most certain revenue.
